\section{Kiến thức nền và công trình liên quan}\label{sec:related}
% Mỗi trích dẫn phải được xác minh (ACL Anthology / arXiv / DBLP) trước khi nộp.

\subsection{MiniRAG}\label{sec:background}
MiniRAG~\cite{fan2025minirag} xây một chỉ mục dị thể gồm hai loại nút: đoạn văn bản và thực thể có tên.
Lúc lập chỉ mục, mô hình ngôn ngữ trích xuất thực thể và quan hệ từ từng đoạn, rồi nối mỗi thực thể với những đoạn chứa nó.
Lúc truy vấn, mô hình phân tích câu hỏi thành từ khoá và loại đáp án dự kiến; hệ thống tìm các thực thể khởi đầu, sinh các đường đi hai bước trên đồ thị và chấm điểm từng đường.
Mỗi đoạn được xếp hạng theo điểm của đường đi chạm tới nó, tạo thành danh sách $L_G$.
Danh sách này được cắt theo ngân sách token và ghép với một bảng thực thể thành ngữ cảnh cho mô hình sinh.
Đối với bài này, chi tiết quan trọng nhất là $L_G$ chỉ chứa những đoạn nằm trên ít nhất một đường đi.
Kết quả truy xuất vectơ chỉ được dùng để cộng điểm cho các đoạn đã có mặt trong $L_G$.

\subsection{RAG dựa trên đồ thị}
Các hệ GraphRAG dùng mô hình ngôn ngữ để dựng đồ thị thực thể từ kho tài liệu, rồi khai thác cấu trúc đồ thị khi truy xuất.
GraphRAG của Edge và cộng sự~\cite{edge2024graphrag} tóm tắt theo cộng đồng thực thể để trả lời câu hỏi mang tính tổng quát; LightRAG~\cite{guo2024lightrag} truy xuất đồng thời ở mức thực thể và mức quan hệ.
MiniRAG đơn giản hoá hướng này để chạy được với mô hình nhỏ.
Bài gốc đánh giá MiniRAG trên LiHua-World, một bộ dữ liệu do chính nhóm tác giả xây dựng, với nhiều SLM \todo{tóm tắt đúng kết luận của bài gốc sau khi đọc lại bản chính thức}.
Các công trình này đánh giá cả hệ thống như một khối; chúng tôi cô lập riêng bước xếp hạng đoạn để đo đóng góp của đồ thị ở đúng bước đó.

\subsection{Truy xuất lai và trộn xếp hạng}
Truy xuất vectơ mã hoá câu hỏi và đoạn văn vào cùng một không gian nhúng~\cite{karpukhin2020dpr,reimers2019sbert}; nó bắt được tương đồng ngữ nghĩa nhưng có thể bỏ sót tên riêng hiếm.
BM25~\cite{robertson2009bm25} xếp hạng theo khớp từ vựng và đến nay vẫn là một đối chứng mạnh.
RRF~\cite{cormack2009rrf} trộn nhiều danh sách chỉ dựa vào thứ hạng, không cần hiệu chuẩn điểm giữa các bộ truy xuất, và thường tốt hơn từng danh sách thành phần.
Truy xuất lai vectơ--BM25 trộn bằng RRF là cấu hình quen thuộc \todo{trích một công trình đánh giá truy xuất lai đã xác minh}.
Chúng tôi không đề xuất cách trộn mới. Các thành phần chuẩn này, với tham số mặc định không tinh chỉnh, đóng vai trò công cụ đo: chúng cho phép hỏi đồ thị có thêm giá trị vào xếp hạng hay không.

\subsection{Đánh giá bằng giám khảo LLM}
Dùng mô hình ngôn ngữ làm giám khảo là cách phổ biến để chấm câu trả lời mở, nhưng giám khảo có thiên lệch riêng và cần được kiểm với người~\cite{zheng2023judge}.
Bài gốc MiniRAG chấm bằng giám khảo LLM với ba nhãn.
Chúng tôi giữ giao thức ba nhãn, thay mô hình giám khảo, và báo kèm độ nhạy theo rubric cùng mức khớp với người chấm.
